
\AddSymbolsF
\pagecolor{ForestGreen!15!Gray!}
\color{white}
\quad\quad \lettrine[,lraise=0.2,lhang=0.5]
{\textbf{教}}{授}和伊藤佐一走进便利店。店内，空调送着冷风，雪柜时常被小孩打开着、偷窥着。

\twkkai{伊藤同学，你认为我们现在的探索处于什么位置？}

教授扶着被推开的便利店的门，突然指着商品柜说。

\twkkai{可能像这一排货柜上的所有商品组成的区域吧；}伊藤佐一支支吾吾地、又略带仓促地应答着。

教授微笑着，但摇了摇头。

\twkkai{那……是这一列吗？}伊藤佐一的信心似乎被打击到了。

教授依旧摇了摇头，接着用手指着地板。

\twkkai{像这一块？}伊藤佐一的语气逐渐犹疑起来。

教授再次摇了摇头，然后终于开口了。\twkkai{不，伊藤同学——我们知道的只是地板上的灰尘，连鞋子都不会理呢。}

伊藤佐一盯着砖缝，悄悄地深吸了一口气，氧气的灌入让他晕乎乎的；他顿觉失落无比。自己所做的一切，在万物看来，似乎都变得无意义了；甚至连铃木和子在这一刻都显得不那么重要了——可是在他心中，她占着多大的比重啊。

教授在雪柜前徘徊着，哼着小调。伊藤佐一则无所适从地尾随着教授，站在他身旁，沉思着。一切都是不确定的——自己未竟的数学探索、未竟的情感追寻——一切，若不在当下立即完成，就仿佛成了未竟带来的遗憾。

可是他才回想起，自己和她，只是两粒灰尘，因风吸引着，又因风排斥着，永不能接触，又永不得安宁。昨天她的忽冷忽热，早已能让他掘地十尺；他无法忘怀——为什么上一秒仍紧密如糖浆，下一秒便一吹即散。

他看着那一格格的砖，以及砖与砖之间的缝隙，发着呆。

教授打断了他的胡思乱想，用一句话将他拉回了现实。

\twkkai{不过啊，伊藤同学，不用感伤——我当年也是这么认为的；后来发现还有比这大得多的东西让我头疼呢。}

他拍了拍伊藤佐一的肩膀。

伊藤佐一愣了一下，接着便看着教授，点了点头。

\twkkai{走吧！哦，对了——你喝饮料吗？}

\twkkai{……不用了，这次就我自己买吧；我还应该请回你才对；}伊藤佐一急忙摆摆手，随即打开了冰柜门。冰凉的饮料刺痛着他的掌心，却仿佛能让他冷静下来。

\twkkai{好啊——谢谢你请我；有没有兴趣去草地上坐坐，聊聊天？}教授打趣地问道。

伊藤佐一的胸腔发热——他有太多太多问题想问教授了；有许多显得像题外话。

……

草地上，两个背影和人们的欢声笑语相对、和天空中的几只风筝呼应着。风筝互相应答；而伊藤佐一也和教授聊得正欢。

\twkkai{教授，其实我一直有个问题想问你；}伊藤佐一清了清喉咙，看着远方的草地，又看向教授的眼镜。

\twkkai{你年轻的时候，也会这样吗——就是，为了一个人，想很多很多事情；}伊藤佐一知道，自己没有后悔问出这句话的余地了。

他没想到，教授几乎像想都没想那样，说了起来。

\twkkai{当然有；我很年轻的时候——当教授之前——深深地喜欢过一个女孩。我对自己说，我一定要和她一起考上数学系——也不知道这句话有没有当着她面说过。我想象过我们的很多未来，比如说考上前怎么一起认真备考、考上后一起逛校园。最终还给她写了一封信；只可惜她没有回。我当时为此伤心了好几天。后来也不知道她的去向了，只有我一个人孤零零地上了大学物理系。}教授摘下了眼镜，揉了揉眼。

\twkkai{原来你是物理系学生？}伊藤佐一眼见如此，假装不经意地问道。

\twkkai{是的；后来，我发现我竟然可以把爱她的动力，转化为我后来调去数学系的动力。}

教授停了下来。而两人头上的两片风筝依然在草地上方左右摆荡着。

\twkkai{那这么说，你一定会后悔什么吧？}伊藤佐一试探着。

\twkkai{那是当然；}教授苦笑了笑，又轻轻地呼出一口气——轻到伊藤佐一几乎听不见。

\twkkai{最后悔的就是为什么自己那么死要面子，不问清楚——我知道自己明明可以亲自去找她的。她不是在教室，就是在图书馆二楼的自习室。要么就是在操场散步——我喜欢跑步，每天都一定会碰上她；只是打不打招呼另说。}

教授又低下了头。伊藤佐一开始感受不到草地的刺痛感，觉得喉咙堵塞。

\twkkai{那教授，如果要你回到那一刻，你会怎么做——或者说，怎么选呢？}

教授推了推眼镜。

\twkkai{要我选的话——其实这几乎不用选——我一定会选择和她在一起，无论结局是考上数学系与否；就像我之前说的，数学太大了，哈哈——大到你追不完。而人心与人心之间是相通的，因此也是可以相互吸引的。}

这时，天上的两只风筝缠在了一起。下方的人们开始呼喊；而教授则轻声地咯咯笑了起来——像个孩子一般。




